\documentclass[12pt,a4paper]{article}

\usepackage[utf8]{inputenc}
\usepackage[T1]{fontenc}
\usepackage[margin=2.5cm]{geometry}
\usepackage{hyperref}
\usepackage{microtype}
\usepackage{parskip}
\usepackage{lmodern}

\hypersetup{colorlinks=true, urlcolor=blue!60!black}

\pagestyle{empty}

\begin{document}

\begin{flushright}
Kundai Farai Sachikonye \\
Auf der Lichtung 19, 82178 Puchheim, Germany \\
\href{https://github.com/fullscreen-triangle}{github.com/fullscreen-triangle} \\
\texttt{kundai.sachikonye@bitspark.com} \\
(+49) 1626386344 \\[6pt]
18 May 2026
\end{flushright}

\vspace{1em}

Prof.\ Dr.\ Anna Poetsch \\
Institut für Genetik / CECAD \\
Universität zu Köln \\
\texttt{apoetsch@uni-koeln.de} \\
Kennziffer: TUV2605-02

\vspace{1.5em}

\textbf{Re: Daten- und IT-Manager*in für Omics-Daten (TUV2605-02)}

\vspace{1em}

Dear Prof.\ Poetsch,

I am applying for the Data and IT Manager position in your group at CECAD
(Kennziffer TUV2605-02).
My background is unusual for this kind of role: I combine hands-on HPC omics
data experience from doctoral research in Computational Lipidomics at TUM
with independent development of production-grade data analysis pipelines
at $>$100\,GB scale and formal work on scientific data management infrastructure.
The combination is directly suited to what the position describes.

\textbf{HPC storage and omics data at scale.}
During my doctoral research at TUM and the Bayerisches Forschungsinstitut
(2019--2021) I worked on mass spectrometry-based lipidomics at HPC scale:
mzML dataset management exceeding 100\,GB per study, federated processing across
distributed clinical sites, peak detection and ion annotation against LIPID MAPS
and LipidBlast, and multi-omics statistical modelling combining lipidomics with
clinical covariates. I am therefore familiar with the practical challenges the
position addresses: data generated faster than it is curated, inconsistent naming
across research groups, and metadata that cannot be recovered without the original
analyst.
The \href{https://github.com/fullscreen-triangle/lavoisier}{Lavoisier} framework
I subsequently built encodes the solutions I worked out to those problems:
a dual-pipeline MS analysis system with memory-mapped I/O for datasets of any
size, structured metadata at ingestion, HDF5 output with consistent schema,
and Ray-distributed processing for multi-node HPC execution.
The 98.9\,\% NIST feature-extraction accuracy is a byproduct of getting the
data structure right first.

\textbf{Metadata schemas, naming conventions, and data lifecycle management.}
The most important practical question in large-scale omics data management is
``given any file on the storage system, can you reconstruct what it contains,
where it came from, and what depends on it?'' --- without asking the person who
created it.
I have formalised this problem mathematically and computationally in the
Purpose-Driven Shader VM (PDSVM) framework, which represents every scientific
data object as a structured coordinate in a three-dimensional metadata space
$(S_k, S_t, S_e) \in [0,1]^3$ (knowledge entropy, temporal entropy, evolution
entropy) and retrieves any object in $O(k)$ operations independent of
collection size, via a ternary trie over Morton-ordered metadata coordinates.
In practice this means: metadata schema design is not a convention problem,
it is a retrieval problem; the schema is correct when any query on the collection
can be answered without full scan.
For the Poetsch group's HPC storage, the same principle translates concretely:
folder hierarchies and naming conventions should be derived from the query patterns
the group actually needs to run (sample, timepoint, treatment, replicate, modality)
rather than from how data happens to arrive; and metadata schemas should be
validated by checking that the $k$ fields jointly address all retrieval cases,
not by committee consensus.

\textbf{SOPs, pipeline standardisation, and ML/DL integration.}
I write SOPs as part of framework development: every tool I publish includes
structured documentation specifying input formats, parameter ranges, expected
outputs, and validation criteria.
For the data analysis pipeline layer, I am experienced with Python and R for
all standard omics operations, and with workflow management in distributed
environments (Ray, Dask, Docker, Kubernetes, FastAPI).
The position mentions ML and Deep Learning in the pipeline context.
The \href{https://github.com/fullscreen-triangle/lavoisier}{Lavoisier} pipeline
includes a CNN-based spectral classification stage (PyTorch) alongside the
numerical pipeline; the \href{https://github.com/fullscreen-triangle/four-sided-triangle}{Four-Sided-Triangle}
RAG system integrates Rust-accelerated (10--50$\times$) inference with
Python orchestration and provides a working template for incorporating
LLM-based document extraction into a scientific data pipeline.
I am therefore positioned to implement DL-augmented analysis steps as
standardised pipeline stages, not as one-off scripts.

\textbf{Bioinformatics ecosystem and genomics data types.}
My MSc in Computational Biology (Universität Tübingen, 2018--2019) covered
the standard NGS stack: sequence alignment, variant calling (VCF), RNA-seq
quantification, and downstream pathway analysis (KEGG, Reactome).
I have practical familiarity with samtools for BAM file operations and with
Nextflow pipeline patterns from the bioinformatics community.
The Poetsch group's data --- including genome-wide sequencing assays relevant
to aging and cancer epigenomics --- is data I understand at both the format
level (FASTQ, BAM, BED, bigWig) and the biological level.
This matters for the part of the role the job description flags explicitly:
consulting researchers on data organisation best practices requires understanding
what their data actually contains and what downstream operations it will need to
support.

\textbf{Communication and interdisciplinary collaboration.}
The position requires working effectively with both technical and non-technical
users. My independent research has required the same skill in a different
register: translating between the mathematical structure of a data problem and
the practical requirements of a wet-lab or clinical collaborator.
I have documented complex frameworks for audiences ranging from algorithm
implementers to experimental biologists, and the resulting documentation is
public and actively maintained.

I hold an MSc in Computational Biology (Universität Tübingen) and an MSc in
Pharmaceutical Biotechnology (HAW Hamburg), and I conducted doctoral research
in Computational Lipidomics at TUM.
I am legally resident in Germany with full work authorisation and am available
from July 2026 or earlier. English is native; German is conversational.

\vspace{1.5em}
Yours sincerely,

\vspace{2.5em}
\textbf{Kundai Farai Sachikonye}

\end{document}
